\chapter{Foundations of Numerical Linear Algebra and SVD Decomposition}
\label{chap:foundations-linear-algebra-svd}

\FloatBarrier
\section{The Role of Linear Algebra in Machine Learning and Data Analysis}

Modern machine learning and large-scale data analysis rely extensively on concepts and methods derived from numerical linear algebra and continuous optimization~\cite{strang2019linear,golub2013matrix}. Real-world data (images, texts, time series, graphs, or sensor measurements) are systematically encoded as numerical vectors and matrices:

\begin{itemize}
    \item \textbf{Design Matrix:} Let $\mathcal{D} = \{\mathbf{x}_1, \dots, \mathbf{x}_n\}$ be a set of $n$ statistical samples, each described by $d$ numerical features. The entire collection is organized as a matrix $X \in \R^{n \times d}$:
    \begin{equation}
        X = \begin{bmatrix}
            \mathbf{x}_1^T \\
            \mathbf{x}_2^T \\
            \vdots \\
            \mathbf{x}_n^T
        \end{bmatrix} = \begin{bmatrix}
            x_{11} & x_{12} & \cdots & x_{1d} \\
            x_{21} & x_{22} & \cdots & x_{2d} \\
            \vdots & \vdots & \ddots & \vdots \\
            x_{n1} & x_{n2} & \cdots & x_{nd}
        \end{bmatrix},
    \end{equation}
    where rows represent observations and columns represent measured variables.
    \item \textbf{Gram Matrix and Covariance Matrix:} Linear relationships between pairs of observations or features are quantified through matrix products: the Gram matrix $G = X X^T \in \R^{n \times n}$ quantifies similarities between data samples, whereas the uncentered sample covariance matrix $\Sigma_X = \frac{1}{n} X^T X \in \R^{d \times d}$ measures variance and pairwise covariance among features.
\end{itemize}

Unlike symbolic or purely theoretical linear algebra, \emph{computational} or \emph{numerical} linear algebra must address two fundamental constraints: asymptotic computational complexity for large-scale operations and error propagation stemming from finite-precision machine arithmetic (IEEE 754 standard)~\cite{trefethen1997numerical}.

\FloatBarrier
\section{Vector and Matrix Norms}

Quantifying the magnitude of vectors and matrices, as well as measuring discrepancies between predictions and ground-truth targets, requires the axiomatic definition of a \textbf{norm}.

\begin{definition}[Vector Norm]
A function $\|\cdot\|: \R^d \to \R$ is called a \textbf{norm} on $\R^d$ if it satisfies the following three properties for all $\mathbf{x}, \mathbf{y} \in \R^d$ and $\alpha \in \R$:
\begin{enumerate}
    \item \textbf{Non-negativity and definiteness:} $\|\mathbf{x}\| \ge 0$, with $\|\mathbf{x}\| = 0 \iff \mathbf{x} = \mathbf{0}$;
    \item \textbf{Absolute homogeneity:} $\|\alpha \mathbf{x}\| = |\alpha| \, \|\mathbf{x}\|$;
    \item \textbf{Triangle inequality:} $\|\mathbf{x} + \mathbf{y}\| \le \|\mathbf{x}\| + \|\mathbf{y}\|$.
\end{enumerate}
\end{definition}

\subsection{The Family of \texorpdfstring{$\ell_p$}{lp}-Norms}

The vector norms most widely encountered in data science belong to the parameterized family of $\ell_p$-norms (with $p \ge 1$):
\begin{equation}
    \|\mathbf{x}\|_p = \left( \sum_{i=1}^d |x_i|^p \right)^{1/p}.
\end{equation}

Three fundamental special cases are:
\begin{enumerate}
    \item \textbf{Manhattan Norm ($\ell_1$):} $\|\mathbf{x}\|_1 = \sum_{i=1}^d |x_i|$. Widely used to promote sparsity in linear models (Lasso regularization).
    \item \textbf{Euclidean Norm ($\ell_2$):} $\|\mathbf{x}\|_2 = \sqrt{\sum_{i=1}^d x_i^2} = \sqrt{\mathbf{x}^T \mathbf{x}}$. Represents standard geometric Euclidean distance and is invariant under orthogonal rotations.
    \item \textbf{Chebyshev Norm ($\ell_\infty$):} $\|\mathbf{x}\|_\infty = \max_{1 \le i \le d} |x_i|$. Captures peak deviation across any individual coordinate.
\end{enumerate}

\begin{property}[Topological Equivalence of Norms]
In a finite-dimensional real vector space $\R^d$, all norms are topologically equivalent: for any pair of norms $\|\cdot\|_a$ and $\|\cdot\|_b$, there exist positive constants $c_1, c_2 > 0$ such that:
\begin{equation}
    c_1 \|\mathbf{x}\|_b \le \|\mathbf{x}\|_a \le c_2 \|\mathbf{x}\|_b \qquad \forall \mathbf{x} \in \R^d.
\end{equation}
\end{property}

\subsection{Matrix Norms and Induced Norms}

To quantify the magnitude of a linear operator $A \in \R^{m \times n}$, matrix norms are employed. Two main classes are distinguished:

\paragraph{Frobenius Norm.}
Directly extends the Euclidean vector norm by treating the matrix as a flattened vector in $\R^{m \cdot n}$:
\begin{equation}
    \|A\|_F = \sqrt{\sum_{i=1}^m \sum_{j=1}^n a_{ij}^2} = \sqrt{\trace(A^T A)}.
\end{equation}

\paragraph{Induced (or Subordinate) Operator Norms.}
Measure the maximum amplification that matrix $A$ can apply to a non-zero vector with respect to a given vector norm:
\begin{equation}
    \|A\|_p = \sup_{\mathbf{x} \ne \mathbf{0}} \frac{\|A \mathbf{x}\|_p}{\|\mathbf{x}\|_p} = \max_{\|\mathbf{x}\|_p = 1} \|A \mathbf{x}\|_p.
\end{equation}

Prominent closed-form computational expressions include:
\begin{itemize}
    \item $\|A\|_1 = \max_{1 \le j \le n} \sum_{i=1}^m |a_{ij}|$ (maximum absolute column sum);
    \item $\|A\|_\infty = \max_{1 \le i \le m} \sum_{j=1}^n |a_{ij}|$ (maximum absolute row sum);
    \item $\|A\|_2 = \sigma_{\max}(A) = \sqrt{\lambda_{\max}(A^T A)}$ (\textbf{spectral norm}, equal to the largest singular value of $A$).
\end{itemize}

Every induced matrix norm is \textbf{sub-multiplicative}: $\|A B\| \le \|A\| \, \|B\|$.

\FloatBarrier
\section{Conditioning and Numerical Stability}

A computational problem $y = f(x)$ is called \textbf{well-conditioned} if small perturbations in the input data $x$ lead to proportionally small variations in the solution $y$. Conversely, a problem is \textbf{ill-conditioned} if slight disturbances in the inputs produce macroscopic fluctuations in the output~\cite{golub2013matrix}.

\subsection{Condition Number of a Linear System}

Consider the square linear system $A \mathbf{x} = \mathbf{b}$, where $A \in \R^{n \times n}$ is non-singular and $\mathbf{b} \ne \mathbf{0}$. Suppose the right-hand side vector is corrupted by additive noise $\delta \mathbf{b}$, inducing a perturbation $\delta \mathbf{x}$ in the solution:
\begin{equation}
    A(\mathbf{x} + \delta \mathbf{x}) = \mathbf{b} + \delta \mathbf{b} \implies A \delta \mathbf{x} = \delta \mathbf{b} \implies \delta \mathbf{x} = A^{-1} \delta \mathbf{b}.
\end{equation}

By the sub-multiplicativity of matrix norms:
\begin{equation}
    \|\delta \mathbf{x}\| \le \|A^{-1}\| \, \|\delta \mathbf{b}\| \qquad \text{and} \qquad \|\mathbf{b}\| \le \|A\| \, \|\mathbf{x}\| \implies \frac{1}{\|\mathbf{x}\|} \le \frac{\|A\|}{\|\mathbf{b}\|}.
\end{equation}

Multiplying both inequalities yields the foundational relative error bound:
\begin{equation}
    \frac{\|\delta \mathbf{x}\|}{\|\mathbf{x}\|} \le \left(\|A\| \, \|A^{-1}\|\right) \frac{\|\delta \mathbf{b}\|}{\|\mathbf{b}\|}.
\end{equation}

\begin{definition}[Condition Number]
The \textbf{condition number} of an invertible square matrix $A \in \R^{n \times n}$ associated with norm $\|\cdot\|$ is:
\begin{equation}
    \kappa(A) = \|A\| \, \|A^{-1}\|.
\end{equation}
In the case of the spectral norm $\|\cdot\|_2$, letting $\sigma_1 \ge \sigma_2 \ge \dots \ge \sigma_n > 0$ denote the singular values of $A$:
\begin{equation}
    \kappa_2(A) = \|A\|_2 \, \|A^{-1}\|_2 = \frac{\sigma_{\max}(A)}{\sigma_{\min}(A)} = \frac{\sigma_1}{\sigma_n} \ge 1.
\end{equation}
\end{definition}

\begin{alertblock}{Computational Significance of $\kappa(A)$}
In standard IEEE 754 double-precision arithmetic (machine precision $\epsilon_{\text{mach}} \approx 10^{-16}$), solving a linear system with $\kappa(A) \approx 10^k$ typically causes the loss of approximately $k$ decimal digits of precision. If $\kappa(A) \ge 10^{16}$, the matrix is regarded as numerically singular.
\end{alertblock}

\FloatBarrier
\section{Spectral Decomposition (Eigendecomposition)}

Let $A \in \R^{n \times n}$ be a square matrix. A scalar $\lambda \in \C$ is an \textbf{eigenvalue} of $A$ with associated non-zero \textbf{eigenvector} $\mathbf{v} \ne \mathbf{0}$ if it satisfies:
\begin{equation}
    A \mathbf{v} = \lambda \mathbf{v}.
\end{equation}

\begin{theorem}[Spectral Theorem for Real Symmetric Matrices]
Let $A \in \R^{n \times n}$ be a real symmetric matrix ($A = A^T$). Then:
\begin{enumerate}
    \item All eigenvalues of $A$ are real: $\lambda_1, \dots, \lambda_n \in \R$.
    \item Eigenvectors corresponding to distinct eigenvalues are mutually orthogonal.
    \item There exists an orthogonal matrix $Q \in \R^{n \times n}$ ($Q^T Q = Q Q^T = I_n$) whose columns are normalized eigenvectors of $A$, such that:
    \begin{equation}
        A = Q \Lambda Q^T = \sum_{i=1}^n \lambda_i \mathbf{q}_i \mathbf{q}_i^T,
    \end{equation}
    where $\Lambda = \diag(\lambda_1, \dots, \lambda_n)$.
\end{enumerate}
\end{theorem}

\subsection{Symmetric Positive Semi-Definite (SPSD) Matrices and Quadratic Forms}

A symmetric matrix $A \in \R^{n \times n}$ is defined as:
\begin{itemize}
    \item \textbf{Positive semi-definite} ($A \succeq 0$) if $\mathbf{x}^T A \mathbf{x} \ge 0$ for all $\mathbf{x} \in \R^n \iff \lambda_i(A) \ge 0$ for all $i=1,\dots,n$.
    \item \textbf{Positive definite} ($A \succ 0$) if $\mathbf{x}^T A \mathbf{x} > 0$ for all $\mathbf{x} \ne \mathbf{0} \iff \lambda_i(A) > 0$ for all $i=1,\dots,n$.
\end{itemize}

\begin{example}[Sample Covariance Matrices are Always Positive Semi-Definite]
Let $X \in \R^{n \times d}$ be a column-centered data matrix. The sample covariance matrix $C = \frac{1}{n} X^T X$ is symmetric and for any arbitrary vector $\mathbf{w} \in \R^d$:
\begin{equation}
    \mathbf{w}^T C \mathbf{w} = \mathbf{w}^T \left(\frac{1}{n} X^T X\right) \mathbf{w} = \frac{1}{n} (X \mathbf{w})^T (X \mathbf{w}) = \frac{1}{n} \|X \mathbf{w}\|_2^2 \ge 0.
\end{equation}
Therefore, all eigenvalues of $C$ are real and non-negative, proving that $C \succeq 0$.
\end{example}

\FloatBarrier
\section{Singular Value Decomposition (SVD)}
\label{sec:svd}

Spectral decomposition applies exclusively to symmetric square matrices. The \textbf{Singular Value Decomposition (SVD)} provides the fundamental generalization to arbitrary rectangular matrices $A \in \R^{m \times n}$~\cite{golub2013matrix,strang2019linear}.

\begin{theorem}[Existence of the SVD]
Given any real matrix $A \in \R^{m \times n}$ of rank $r \le \min(m, n)$, there exist two orthogonal matrices $U \in \R^{m \times m}$ and $V \in \R^{n \times n}$ and a pseudo-diagonal matrix $\Sigma \in \R^{m \times n}$ such that:
\begin{equation}
    A = U \Sigma V^T,
\end{equation}
where the main diagonal of $\Sigma$ contains the singular values arranged in non-increasing order:
\begin{equation}
    \sigma_1 \ge \sigma_2 \ge \dots \ge \sigma_r > \sigma_{r+1} = \dots = \sigma_{\min(m, n)} = 0.
\end{equation}
\end{theorem}

\begin{figure}[H]
\centering
\begin{tikzpicture}[
    boxU/.style={rectangle, draw=blue!80!black, fill=blue!10, thick, minimum width=1.8cm, minimum height=2.8cm, align=center, font=\small\bfseries},
    boxS/.style={rectangle, draw=green!70!black, fill=green!10, thick, minimum width=2.4cm, minimum height=2.8cm, align=center, font=\small\bfseries},
    boxV/.style={rectangle, draw=orange!80!black, fill=orange!10, thick, minimum width=2.4cm, minimum height=1.6cm, align=center, font=\small\bfseries},
    boxA/.style={rectangle, draw=purple!80!black, fill=purple!10, thick, minimum width=2.4cm, minimum height=2.8cm, align=center, font=\small\bfseries},
    lbl/.style={font=\small}
]
    \node[boxA] (A) at (0, 0) {$A$\\$(m \times n)$};
    \node[lbl] at (1.8, 0) {$=$};
    \node[boxU] (U) at (3.2, 0) {$U$\\$(m \times m)$};
    \node[boxS] (S) at (5.7, 0) {$\Sigma$\\$(m \times n)$};
    \node[boxV] (V) at (8.5, 0.6) {$V^T$\\$(n \times n)$};

    \node[lbl, below=0.2cm of A, text=purple!80!black] {Original data};
    \node[lbl, below=0.2cm of U, text=blue!80!black] {Left singular vectors};
    \node[lbl, below=0.2cm of S, text=green!70!black] {Singular values};
    \node[lbl, below=1.4cm of V, text=orange!80!black] {Right singular vectors};
\end{tikzpicture}
\caption{Dimensional block representation of the Singular Value Decomposition (Full SVD).}
\label{fig:svd-blocks-en}
\end{figure}

\subsection{Outer Product Expansion and Thin SVD}

The matrix $A$ can be decomposed into an outer-product sum of $r$ rank-one matrices:
\begin{equation}
    A = \sum_{i=1}^r \sigma_i \mathbf{u}_i \mathbf{v}_i^T,
\end{equation}
where:
\begin{itemize}
    \item $\mathbf{u}_i \in \R^m$ is the $i$-th column of $U$ (\textbf{left singular vector}), an eigenvector of $A A^T$ with eigenvalue $\lambda_i = \sigma_i^2$;
    \item $\mathbf{v}_i \in \R^n$ is the $i$-th column of $V$ (\textbf{right singular vector}), an eigenvector of $A^T A$ with eigenvalue $\lambda_i = \sigma_i^2$;
    \item $\sigma_i = \sqrt{\lambda_i(A^T A)} = \sqrt{\lambda_i(A A^T)}$ is the $i$-th singular value.
\end{itemize}

In computational settings where $m \gg n$, the \textbf{thin SVD} (or economy SVD) is used instead: it retains only the first $n$ columns of $U$ and the square diagonal block $\Sigma_n \in \R^{n \times n}$, reducing memory footprint from $\BigO{m^2}$ down to $\BigO{m n}$.

\FloatBarrier
\section{Properties and Applications of the SVD}

\subsection{Best Low-Rank Approximation: Eckart-Young-Mirsky Theorem}

In exploratory data analysis, signal denoising, and lossy compression, one frequently seeks a low-rank matrix $\hat{A}_k$ with rank $k < r$ that best approximates the original matrix $A$.

\begin{theorem}[Eckart-Young-Mirsky Theorem]
Let $A = \sum_{i=1}^r \sigma_i \mathbf{u}_i \mathbf{v}_i^T$ be the SVD of $A \in \R^{m \times n}$. For any integer $1 \le k < r$, the truncated SVD matrix:
\begin{equation}
    A_k = \sum_{i=1}^k \sigma_i \mathbf{u}_i \mathbf{v}_i^T = U_k \Sigma_k V_k^T
\end{equation}
is the unique optimal solution to the low-rank approximation problem in both the spectral norm and the Frobenius norm:
\begin{equation}
    A_k = \argmin_{B \in \R^{m \times n},\, \rank(B) \le k} \|A - B\|_2 = \argmin_{B \in \R^{m \times n},\, \rank(B) \le k} \|A - B\|_F.
\end{equation}
Moreover, the optimal approximation errors are given explicitly by:
\begin{align}
    \|A - A_k\|_2 &= \sigma_{k+1}, \\
    \|A - A_k\|_F &= \sqrt{\sum_{j=k+1}^r \sigma_j^2}.
\end{align}
\end{theorem}

\subsection{The Moore-Penrose Pseudoinverse and Linear Least Squares}

When a rectangular matrix $A \in \R^{m \times n}$ lacks a standard inverse ($m \ne n$ or rank-deficient), the SVD yields the canonical formulation of the \textbf{Moore-Penrose pseudoinverse} $A^\dagger \in \R^{n \times m}$:
\begin{equation}
    A^\dagger = V \Sigma^\dagger U^T = \sum_{i=1}^r \frac{1}{\sigma_i} \mathbf{v}_i \mathbf{u}_i^T,
\end{equation}
where $\Sigma^\dagger$ is obtained by transposing $\Sigma$ and inverting only the strictly positive diagonal entries $\sigma_i > 0$.

\begin{property}[Least Squares Property]
For the linear least squares problem $\min_{\mathbf{x} \in \R^n} \|A \mathbf{x} - \mathbf{b}\|_2$, the vector
\begin{equation}
    \mathbf{x}^* = A^\dagger \mathbf{b}
\end{equation}
is the unique minimum-norm solution $\|\mathbf{x}^*\|_2$ among all vectors minimizing the residual sum of squares.
\end{property}

\subsection{Connection Between SVD and Principal Component Analysis (PCA)}

Let $X \in \R^{n \times d}$ denote a column-centered data matrix (empirical sample mean $\boldsymbol{\mu} = \mathbf{0}$). The empirical covariance matrix is:
\begin{equation}
    C = \frac{1}{n-1} X^T X.
\end{equation}
Expressing the SVD of $X$ as $X = U \Sigma V^T$, we obtain:
\begin{equation}
    C = \frac{1}{n-1} (V \Sigma U^T) (U \Sigma V^T) = \frac{1}{n-1} V \Sigma^2 V^T = V \left( \frac{\Sigma^2}{n-1} \right) V^T.
\end{equation}

From this identity, three crucial insights follow directly:
\begin{enumerate}
    \item The right singular vectors $\mathbf{v}_i$ of $X$ (columns of $V$) are precisely the \textbf{eigenvectors of the sample covariance matrix} (PCA principal directions).
    \item The variance explained by the $i$-th principal component is directly proportional to the squared singular value: $\lambda_i = \frac{\sigma_i^2}{n-1}$.
    \item Low-dimensional data projections (\textit{principal component scores}) onto the top $k$ components can be obtained directly without ever instantiating $C \in \R^{d \times d}$:
    \begin{equation}
        Z_k = X V_k = (U_k \Sigma_k V_k^T) V_k = U_k \Sigma_k.
    \end{equation}
\end{enumerate}

\begin{lstlisting}[style=pythoncode,caption={Compact Truncated SVD and PCA computation using NumPy.}]
import numpy as np

def compute_svd_pca(X: np.ndarray, k: int):
    """
    Computes a k-dimensional PCA projection via the economy SVD.
    X: data matrix of shape (n_samples, n_features)
    k: number of principal components to retain
    """
    # 1. Center the data (zero-mean along columns)
    X_centered = X - np.mean(X, axis=0)
    n = X.shape[0]

    # 2. Compute economy SVD
    U, s, Vt = np.linalg.svd(X_centered, full_matrices=False)

    # 3. Truncate to top-k components
    Uk = U[:, :k]
    sk = s[:k]
    Vk = Vt[:k, :].T

    # 4. Compute projected coordinates (scores) and explained variance
    scores = Uk * sk  # equivalent to X_centered @ Vk
    explained_variance = (sk ** 2) / (n - 1)
    explained_variance_ratio = explained_variance / np.sum((s ** 2) / (n - 1))

    return scores, Vk, explained_variance_ratio
\end{lstlisting}

\FloatBarrier
\section{Summary of Core Mathematical Concepts}

\begin{table}[H]
  \centering
  \small
  \setlength{\tabcolsep}{6pt}
  \begin{tabularx}{\textwidth}{l Y Y}
    \toprule
    \textbf{Mathematical Tool} & \textbf{Operational Definition} & \textbf{Impact in Data Analysis / ML} \\
    \midrule
    $\ell_2$ norm and $\|A\|_F$ & Square root of sum of squared entries & Least squares optimization, MSE risk metric. \\
    $\ell_1$ norm & Sum of absolute values & Sparsity induction, robust feature selection (Lasso). \\
    Condition number $\kappa$ & $\sigma_{\max} / \sigma_{\min}$ & Numerical stability, noise amplification in data. \\
    Eigendecomposition & $A = Q \Lambda Q^T$ (only for $A = A^T$) & Spectral analysis, quadratic forms, positive kernels. \\
    Full SVD & $A = U \Sigma V^T$ (arbitrary rectangular) & Universal geometric decomposition of data matrices. \\
    Eckart-Young-Mirsky & $A_k = \sum_{i=1}^k \sigma_i \mathbf{u}_i \mathbf{v}_i^T$ & Optimal low-rank compression, noise filtering. \\
    SVD for PCA & $Z_k = U_k \Sigma_k$ & Dimension reduction without forming $X^T X$. \\
    \bottomrule
  \end{tabularx}
  \caption{Summary of numerical linear algebra and SVD concepts covered in this chapter.}
  \label{tab:summary-linear-algebra-svd-en}
\end{table}
